% !TEX root = ../../tcc.tex

\chapter[Referencial Teórico]{Referencial Teórico}
\label{cap:referencial}

A compreensão do sistema proposto demanda uma base conceitual sólida acerca dos tópicos abordados ao longo do trabalho, contemplando a engenharia de software, em suas dimensões de engenharia de requisitos orientada à \textit{Lean Inception}, de metodologias ágeis de desenvolvimento e de arquitetura de sistemas intensivos em dados, além do \textit{Learning Analytics} e da aprendizagem de máquina aplicada ao rastreamento de conhecimento (\textit{Knowledge Tracing}).

\section{Engenharia de Software}
\label{sec:engenharia-software}

A Engenharia de Software é uma disciplina da engenharia que compreende todos os aspectos da produção de sistemas de software, desde os estágios iniciais de especificação e concepção até a operação, manutenção e evolução contínua após a implantação \cite{sommerville_engenharia_2019}. Diferentemente da programação isolada de algoritmos, a engenharia de software preocupa-se com a aplicação de abordagens sistemáticas, disciplinadas e quantificáveis ao desenvolvimento, à operação e à manutenção de software, buscando assegurar atributos de qualidade como confiabilidade, manutenibilidade, segurança e eficiência sob restrições reais de prazo, custo e governança \cite{bourque_guide_2014, pressman_engenharia_2021}.

Segundo \citeonline{sommerville_engenharia_2019}, todo processo de software compreende quatro atividades fundamentais: especificação, projeto e implementação, validação e evolução. Este capítulo aborda as que mais se relacionam ao planejamento do módulo Cognita: a especificação, por meio da engenharia de requisitos, e o projeto, por meio da arquitetura de software, além das metodologias que organizam o desenvolvimento.

Em plataformas educacionais distribuídas e orientadas a dados, como o Ecossistema CEDIS, a Engenharia de Software precisa articular componentes transacionais com fluxos analíticos e modelos de aprendizagem de máquina \cite{kleppmann_designing_2017}. Assim, a construção de um módulo que estima o domínio de conceitos dos estudantes exige, além do modelo de \textit{Knowledge Tracing}, cuidado na elicitação e na priorização dos requisitos, no desenvolvimento incremental do produto e no projeto de uma arquitetura desacoplada, que preserve a integridade dos dados, a reprodutibilidade das análises e a privacidade dos estudantes.

\subsection{Engenharia de Requisitos}
\label{subsec:engenharia-requisitos}

A Engenharia de Requisitos (ER) consiste no processo sistemático de descobrir, analisar, documentar, validar e gerenciar os serviços que um sistema deve fornecer e as restrições operacionais e de domínio sob as quais ele deve funcionar \cite{sommerville_engenharia_2019, wiegers_software_2013}. Falhas, ambiguidades ou desalinhamentos nessa etapa figuram historicamente entre as principais causas de retrabalho e construção de funcionalidades sem valor real em projetos de software \cite{vazquez_engenharia_2016}. Na literatura clássica, os requisitos de software dividem-se em duas categorias complementares \cite{sommerville_engenharia_2019, pressman_engenharia_2021}:

\begin{itemize}
    \item \textbf{Requisitos Funcionais (RF):} descrevem explicitamente o que o sistema deve fazer, detalhando os serviços e funcionalidades entregues aos usuários, as entradas processadas e as saídas geradas. No escopo deste trabalho, os requisitos funcionais abrangem a reconciliação entre respostas de atividades e conceitos do domínio, a estimativa do domínio de conceitos de cada estudante e a apresentação de diagnósticos acionáveis para estudantes e docentes;
    \item \textbf{Requisitos Não Funcionais (RNF):} especificam restrições e atributos de qualidade emergentes do sistema, como desempenho, interoperabilidade, usabilidade, segurança e privacidade, os quais condicionam diretamente as decisões arquiteturais do software \cite{bass_software_2021}.
\end{itemize}

Em plataformas educacionais que coletam registros de interação de estudantes, assumem protagonismo os requisitos não funcionais de proteção de dados e privacidade desde a concepção (\textit{Privacy by Design}). Formulado por \citeonline{cavoukian_privacy_2009}, o paradigma de \textit{Privacy by Design} estabelece que a privacidade deve ser incorporada preventivamente à arquitetura do sistema como configuração padrão (\textit{privacy as the default setting}), e não como ajuste posterior. No contexto brasileiro, esse requisito traduz-se na conformidade estrita com a Lei Geral de Proteção de Dados (LGPD), impondo minimização de coleta, pseudonimização de telemetria e segregação entre dados operacionais anônimos e identidades fornecidas mediante consentimento explícito (\textit{opt-in}).

Em projetos orientados por métodos ágeis e enxutos, a Engenharia de Requisitos afasta-se da produção de extensas especificações estáticas preliminares e adota uma abordagem colaborativa, iterativa e centrada no valor entregue ao usuário \cite{vazquez_engenharia_2016, cohn_user_2004}. Em vez de detalhar exaustivamente todas as funções de antemão, o levantamento e a priorização de requisitos concentram-se na descoberta contínua do produto, combinando técnicas visuais de alinhamento de escopo, como a metodologia \textit{Lean Inception} \cite{caroli_lean_2018}, com a prototipação de alta fidelidade \cite{pressman_engenharia_2021}, que tangibiliza as interfaces analíticas para validação prévia junto aos \textit{stakeholders}.

\subsubsection{\textit{Lean Inception}}
\label{subsubsec:lean-inception}

Proposta por \citeonline{caroli_lean_2018}, a \textit{Lean Inception} é uma metodologia colaborativa projetada para alinhar um grupo de pessoas em torno da descoberta, elicitação e priorização dos requisitos essenciais para construir o produto certo, culminando na definição do Produto Mínimo Viável (\textit{Minimum Viable Product}, MVP). A metodologia nasce da convergência direta entre duas abordagens fundamentais da engenharia de produtos digitais: o \textit{Design Thinking} \cite{brown_design_2010}, que coloca a empatia com as necessidades humanas, a ideação colaborativa e a experimentação visual no centro da resolução de problemas, e o \textit{Lean Startup} \cite{ries_lean_2011}, que preconiza a eliminação de desperdícios e a validação rápida de hipóteses de valor por meio do ciclo construir-medir-aprender (\textit{build-measure-learn}).

Dentro do conceito de \citeonline{ries_lean_2011}, o MVP não corresponde a um produto incompleto ou desprovido de qualidade técnica, mas sim à menor interseção funcional capaz de entregar a proposta de valor central aos usuários, validar premissas críticas do domínio e gerar aprendizado validado com o mínimo esforço de engenharia. Para transformar objetivos abstratos em um conjunto estruturado e priorizado de requisitos do MVP, a \textit{Lean Inception} organiza o processo de engenharia de requisitos em uma sequência lógica de dinâmicas e artefatos \cite{caroli_lean_2018}:

\begin{enumerate}
    \item \textbf{Visão do Produto:} estabelece de forma concisa e compartilhada o propósito da solução, o público-alvo, a necessidade ou problema endereçado, o nome do produto/módulo e seu diferencial em relação às alternativas existentes;
    \item \textbf{É/Não é/Faz/Não faz:} delimita explicitamente as fronteiras de escopo do sistema, prevenindo ambiguidades ao explicitar tanto o que o software realiza quanto o que está fora de seu escopo;
    \item \textbf{Objetivos do Negócio:} alinha as metas estratégicas esperadas pelas partes interessadas (como aprimorar o diagnóstico pedagógico e reduzir a evasão ou defasagem conceitual);
    \item \textbf{Personas:} modela perfis representativos dos usuários reais do sistema (no contexto deste trabalho, os estudantes e os docentes do Ecossistema CEDIS), descrevendo seus comportamentos, dores e necessidades informacionais;
    \item \textbf{Jornadas do Usuário:} mapeia o fluxo sequencial de passos que cada persona percorre antes, durante e após interagir com o ecossistema, revelando os pontos de contato onde o sistema deve intervir;
    \item \textbf{\textit{Brainstorming} de Funcionalidades:} traduz os objetivos de negócio e os pontos de dor das jornadas em requisitos funcionais candidatos (funcionalidades), descrevendo ações ou interações concretas que geram valor para as personas;
    \item \textbf{Revisão Técnica, de Negócio e de UX:} submete cada funcionalidade candidata a uma avaliação multidimensional de nível de confiança (técnico e arquitetural) e a uma matriz de esforço (\textit{E}), valor de negócio (\$) e valor de experiência do usuário ($\heartsuit$), permitindo a priorização objetiva dos requisitos;
    \item \textbf{Sequenciador:} organiza as funcionalidades priorizadas em ondas incrementais de entrega, respeitando dependências técnicas e limites máximos de esforço por onda, de modo a demarcar com precisão quais funcionalidades compõem o MVP e quais constituem incrementos subsequentes;
    \item \textbf{\textit{Canvas} MVP:} consolida, em um único quadro visual de engenharia de requisitos, as personas segmentadas, a proposta de valor, as funcionalidades do MVP, o custo e cronograma previstos, os resultados esperados e as métricas de validação das hipóteses do produto.
\end{enumerate}

Neste trabalho, a adoção de uma versão adaptada da \textit{Lean Inception} permitiu conduzir a Engenharia de Requisitos de maneira enxuta e rastreável: partindo das restrições reais da arquitetura de dados do Ecossistema CEDIS e das demandas pedagógicas de estudantes e professores, as funcionalidades do módulo de estimativa do domínio de conceitos foram elicitadas, avaliadas quanto à viabilidade técnica e ao valor analítico, e organizadas em um MVP coeso para implementação incremental.

\subsection{Metodologias de Desenvolvimento de Software}
\label{subsec:metodologias-desenvolvimento}

As metodologias de desenvolvimento de software estruturam o ciclo de vida de construção do sistema por meio de práticas, papéis, cerimônias e artefatos que orientam a transformação dos requisitos especificados em software operacional \cite{pressman_engenharia_2021}. Em contraste com os modelos tradicionais orientados a planos rígidos, como o modelo em cascata (\textit{Waterfall}), no qual todas as fases são executadas de maneira estritamente sequencial e as mudanças têm alto custo, as metodologias ágeis consolidaram-se para responder a cenários de evolução contínua, complexidade técnica e necessidade de entregas frequentes de valor \cite{sommerville_engenharia_2019}.

Fundamentado nos valores e princípios do Manifesto Ágil \cite{beck_manifesto_2001}, o desenvolvimento ágil prioriza indivíduos e interações, software em funcionamento, colaboração com as partes interessadas e capacidade de resposta a mudanças através de ciclos iterativos e incrementais. Enquanto a \textit{Lean Inception} atua na fase de descoberta e alinhamento inicial para definir o que deve ser construído no MVP, os \textit{frameworks} ágeis de execução, com destaque para o Scrum \cite{schwaber_scrum_2020}, sistematizam como esse produto é construído e evoluído ao longo do tempo.

O Scrum baseia-se no empirismo (transparência, inspeção e adaptação) e organiza o trabalho em ciclos iterativos de duração fixa denominadas \textit{Sprints} \cite{schwaber_scrum_2020}. As funcionalidades e requisitos priorizados no Sequenciador e no \textit{Canvas} MVP da \textit{Lean Inception} alimentam diretamente o \textit{Backlog} do Produto (\textit{Product Backlog}), sendo refinados em itens acionáveis e histórias de usuário \cite{cohn_user_2004, caroli_lean_2018}. No início de cada \textit{Sprint}, realiza-se o planejamento (\textit{Sprint Planning}) para selecionar o subconjunto de itens que comporá o \textit{Backlog} da \textit{Sprint}, estabelecendo uma meta clara de entrega. Ao final de cada ciclo, o incremento funcional produzido passa por revisão (\textit{Sprint Review}) e o processo de trabalho é avaliado em retrospectiva (\textit{Sprint Retrospective}), garantindo previsibilidade no cronograma de desenvolvimento e mitigação contínua de riscos técnicos.

\subsection{Design e Arquitetura de Software}
\label{subsec:design-arquitetura}

Enquanto o design de software em nível detalhado ocupa-se da estrutura interna de módulos, classes, funções e interfaces de usuário, a arquitetura de software representa a estrutura de mais alto nível de um sistema, compreendendo os elementos de software, suas propriedades externamente visíveis e os relacionamentos que os conectam \cite{bass_software_2021}. De acordo com \citeonline{richards_fundamentals_2020}, a arquitetura de software é definida pela combinação de quatro dimensões: a estrutura arquitetural adotada (como estilos em camadas, microsserviços ou orientados a eventos), as características arquiteturais (atributos de qualidade como escalabilidade, segurança, auditabilidade e disponibilidade), as decisões arquiteturais (regras e restrições de projeto) e os princípios de design que guiam a implementação.

Em ecossistemas educacionais compostos por múltiplos produtos independentes, como o Ecossistema CEDIS, adota-se predominantemente uma arquitetura distribuída orientada a serviços, na qual cada componente possui responsabilidades e domínios de dados estritamente delimitados \cite{richards_fundamentals_2020}. Nesse cenário de desacoplamento, a comunicação intercomponencial realiza-se por meio de Interfaces de Programação de Aplicações (APIs) regidas por contratos explícitos de dados e protocolos de segurança padronizados, tais como OAuth2 e \textit{OpenID Connect} (OIDC) para autenticação e autorização, e códigos de autenticação de mensagem baseados em \textit{hash} (HMAC) para verificação de integridade e autenticidade entre serviços.

Uma consequência arquitetural fundamental de sistemas distribuídos com bancos de dados independentes por serviço é a substituição de chaves estrangeiras físicas por referências lógicas entre domínios \cite{kleppmann_designing_2017}. Diferentemente de arquiteturas monolíticas apoiadas em um único sistema gerenciador de banco de dados relacional, a integridade referencial e a consistência dos dados entre múltiplos serviços dependem de esquemas bem definidos (\textit{schemas}), validação contratual nas bordas de ingestão e mecanismos de reconciliação assíncrona.

Quando o sistema precisa consolidar registros operacionais heterogêneos para alimentar modelos analíticos e de aprendizagem de máquina, a arquitetura deve incorporar padrões específicos de aplicações intensivas em dados (\textit{data-intensive applications}) \cite{kleppmann_designing_2017}. Nesses sistemas, separa-se o processamento transacional online (\textit{Online Transaction Processing}, OLTP), otimizado para baixa latência no registro de interações e respostas dos usuários, do processamento analítico (\textit{Online Analytical Processing}, OLAP) e dos \textit{pipelines} de engenharia de dados. Conforme sintetizado na Tabela~\ref{tab:conceitos-engenharia-software}, a articulação entre a engenharia de requisitos guiada pela \textit{Lean Inception}, o desenvolvimento ágil em \textit{Sprints} e a arquitetura de sistemas intensivos em dados fornece o alicerce estrutural para a concepção do módulo de \textit{Knowledge Tracing} proposto neste trabalho.

\begin{table}[H]
\caption{Síntese dos conceitos de Engenharia de Software aplicados ao projeto}
\label{tab:conceitos-engenharia-software}
\centering
\begin{tabularx}{\textwidth}{l>{\raggedright\arraybackslash}X>{\raggedright\arraybackslash}X}
\hline
\textbf{Subárea} & \textbf{Conceitos e práticas centrais} & \textbf{Aplicação direta no TCC} \\
\hline
Engenharia de Requisitos e \textit{Lean Inception} & Requisitos Funcionais (RF), Requisitos Não Funcionais (RNF), \textit{Privacy by Design}, \textit{Lean Inception}, MVP e prototipação de alta fidelidade \cite{sommerville_engenharia_2019, cavoukian_privacy_2009, ries_lean_2011, caroli_lean_2018} & Alinhamento de visão, personas e jornadas; priorização das funcionalidades do MVP de diagnóstico conceitual; conformidade \textit{LGPD-first} e validação visual das interfaces \\
\hline
Metodologias de Desenvolvimento & Desenvolvimento iterativo e incremental, Manifesto Ágil, Scrum (\textit{Product Backlog} e \textit{Sprints}) \cite{beck_manifesto_2001, schwaber_scrum_2020, cohn_user_2004} & Desdobramento do Sequenciador da \textit{Lean Inception} em \textit{Backlog} e estruturação do cronograma incremental de desenvolvimento em \textit{Sprints} \\
\hline
Design e Arquitetura de Software & Sistemas distribuídos, integração via contratos de API, referências lógicas e \textit{pipelines} intensivos em dados \cite{bass_software_2021, richards_fundamentals_2020, kleppmann_designing_2017} & Reconciliação de múltiplas fontes do Ecossistema CEDIS (Slides, Deck Manager, Atlas e Learning Passport) em um \textit{pipeline} de atributos para \textit{Knowledge Tracing} \\
\hline
\end{tabularx}
\legend{Fonte: Elaborado pelos autores.}
\end{table}




\section{Learning Analytics}
\label{sec:learning-analytics}

O campo de \textit{Learning Analytics} (LA) consolidou-se na interseção entre a ciência da computação, a estatística aplicada, as ciências cognitivas e a educação como resposta à crescente disponibilidade de rastros digitais gerados por ambientes virtuais e ecossistemas tecnológicos de ensino \cite{lang_handbook_2017}. Em sua definição canônica, estabelecida na primeira conferência internacional da área (\textit{International Conference on Learning Analytics and Knowledge}, LAK) e adotada pela \textit{Society for Learning Analytics Research} (SoLAR), \textit{Learning Analytics} consiste na ``medição, coleta, análise e comunicação de dados sobre os aprendizes e seus contextos, com o propósito de compreender e otimizar a aprendizagem e os ambientes nos quais ela ocorre'' \cite{siemens_learning_2012, lang_handbook_2017}.

Na literatura especializada, o \textit{Learning Analytics} desenvolveu-se em estreito diálogo com a Mineração de Dados Educacionais (\textit{Educational Data Mining}, EDM), compartilhando técnicas de modelagem estatística e aprendizagem de máquina, mas distinguindo-se em sua ênfase sistêmica e humana \cite{siemens_learning_2012}. Enquanto a EDM historicamente priorizou a descoberta algorítmica de padrões e a adaptação automatizada dentro de tutores computacionais, incluindo a formulação matemática de modelos de \textit{Knowledge Tracing}, o \textit{Learning Analytics} enfatiza a interpretação contextualizada do fenômeno educativo, a transparência dos indicadores e o empoderamento da tomada de decisão humana por parte de estudantes e professores (\textit{human-in-the-loop}) \cite{siemens_learning_2012, gaevic_learning_2015}. No contexto do módulo Cognita proposto neste trabalho, essas duas vertentes convergem de maneira complementar: a modelagem probabilística de rastreamento de conhecimento fornece o motor preditivo de estimação de domínio, ao passo que os princípios de \textit{Learning Analytics} estruturam o ciclo de coleta multi-fonte, a reconciliação semântica, a explicabilidade das estimativas, a governança ética de privacidade e a entrega de diagnósticos acionáveis aos atores do processo pedagógico.

Para sistematizar como dados brutos de interação convertem-se em ganho educacional efetivo, \citeonline{clow_learning_2012} formulou o \textbf{Ciclo de \textit{Learning Analytics}}, estruturado em quatro estágios interdependentes:
\begin{enumerate}
    \item \textbf{Aprendizes (\textit{Learners}):} estudantes engajados em situações formais ou autônomas de aprendizagem, como a participação em aulas com atividades interativas ou a resolução de questões em ambientes de prática;
    \item \textbf{Dados (\textit{Data}):} captura e armazenamento de rastros comportamentais, submissões de respostas avaliativas, carimbos temporais e metadados de contexto;
    \item \textbf{Métricas e Análise (\textit{Metrics / Analytics}):} processamento, validação, engenharia de atributos e aplicação de modelos descritivos, diagnósticos ou preditivos que transformam registros transacionais em indicadores sobre o processo de aprendizagem;
    \item \textbf{Intervenções (\textit{Interventions}):} fechamento do ciclo (\textit{closing the loop}) por meio do retorno de \textit{feedback}, recomendações ou alertas que alteram a ação do estudante ou o planejamento didático do professor.
\end{enumerate}

Segundo \citeonline{clow_learning_2012} e \citeonline{gaevic_learning_2015}, uma falha recorrente em plataformas educacionais consiste em interromper esse ciclo no terceiro estágio, acumulando painéis de telemetria operacional que não se traduzem em ação pedagógica fundamentada. Superar essa limitação exige compreender os diferentes níveis de maturidade analítica e os requisitos arquiteturais, cognitivos e éticos que regem sistemas modernos de \textit{Learning Analytics}.

\subsection{Maturidade Analítica: da Telemetria Descritiva ao Diagnóstico Cognitivo}
\label{subsec:maturidade-la}

A evolução da capacidade analítica em plataformas educacionais é caracterizada na literatura por uma progressão em quatro níveis de maturidade: descritivo, diagnóstico, preditivo e prescritivo/orientador \cite{lang_handbook_2017, liao_naviedu_2026}. No estágio inicial, a \textbf{análise descritiva}, os sistemas concentram-se em resumir eventos passados de interação, reportando métricas de volume e fluxo, tais como número de acessos a materiais, impressões de tela, taxas de início e submissão de atividades, e distribuições percentuais de respostas. Embora essenciais para monitorar a operação técnica e o engajamento superficial, indicadores puramente descritivos apresentam uma limitação epistemológica central para a educação: \textbf{engajamento comportamental e exposição a conteúdo não equivalem a aprendizagem ou domínio conceitual} \cite{gaevic_learning_2015}. Um estudante pode abrir repetidamente uma apresentação de slides ou submeter respostas em sala de aula sem ter compreendido o conceito subjacente, assim como médias gerais de participação de uma turma podem ocultar lacunas severas em pré-requisitos estruturais do currículo.

A transição para a \textbf{análise diagnóstica} ocorre quando os registros de interação deixam de ser tratados como eventos isolados de interface e passam a ser reconciliados com a estrutura semântica do domínio de conhecimento, tipicamente representada por mapas conceituais e por uma Matriz Q que associa cada questão ou atividade avaliativa aos conceitos de conhecimento (\textit{Knowledge Concepts}, KC) exigidos para sua resolução \cite{cao_multi-granularity_2025, liu_enhancing_2026}. Nesse nível, o sistema torna-se capaz de isolar quais conceitos específicos explicam as dificuldades observadas nas respostas de um estudante ou de uma coorte.

O terceiro nível, a \textbf{análise preditiva}, incorpora a dimensão temporal e probabilística ao diagnóstico. Reconhecendo que o conhecimento humano é um estado latente (não diretamente observável) e dinâmico, sujeito a acertos casuais por chute (\textit{guess}) e a erros acidentais por descuido (\textit{slip}), modelos de \textit{Knowledge Tracing} estimam a probabilidade atualizada de domínio $P(L_t)$ de cada conceito ao longo de uma sequência cronológica de tentativas $X = \{(q_1, a_1), \ldots, (q_t, a_t)\}$, permitindo prever o desempenho em interações futuras \cite{corbett_knowledge_1994, pelanek_bayesian_2017}.

Por fim, o nível \textbf{prescritivo ou orientador} fecha o ciclo proposto por \citeonline{clow_learning_2012} ao converter as probabilidades inferidas e os diagnósticos de lacunas em suporte direto à decisão: para o estudante, na forma de priorização explícita de quais conceitos revisar primeiro; para o docente, na sinalização dos conceitos com baixa absorção coletiva e dos materiais didáticos que demandam reforço \cite{liao_naviedu_2026}. A Tabela~\ref{tab:niveis-learning-analytics} sintetiza esses quatro níveis e explicita como o módulo Cognita promove a evolução analítica do Ecossistema CEDIS.

\begin{table}[H]
\caption{Níveis de maturidade em \textit{Learning Analytics} e sua correspondência no Ecossistema CEDIS e no módulo Cognita}
\label{tab:niveis-learning-analytics}
\centering
\begin{tabularx}{\textwidth}{>{\raggedright\arraybackslash}p{2.7cm}>{\raggedright\arraybackslash}p{2.8cm}>{\raggedright\arraybackslash}X>{\raggedright\arraybackslash}X}
\hline
\textbf{Nível Analítico} & \textbf{Questão Central} & \textbf{Técnicas e Artefatos Típicos} & \textbf{Materialização no CEDIS / Cognita} \\
\hline
\textbf{Descritivo} \cite{siemens_learning_2012, lang_handbook_2017} & O que aconteceu? & Contagens de eventos, funis de navegação, visitantes por dia (\textit{visitor-days}) e distribuição bruta de respostas. & Telemetria operacional já existente no Deck Manager (\textit{DeckEvent}) e resumos de participação do CEDIS Slides. \\
\hline
\textbf{Diagnóstico} \cite{gaevic_learning_2015, cao_multi-granularity_2025} & Onde e por que ocorreu a lacuna de aprendizagem? & Reconciliação entre respostas e estrutura semântica do domínio (mapas conceituais e Matriz Q); isolamento de conceitos deficitários. & Cruzamento de respostas (\textit{activity\_responses}) com a Matriz Q do Concept Studio + Questionários para identificar conceitos com baixa absorção e explicar seus motivos. \\
\hline
\textbf{Preditivo} \cite{corbett_knowledge_1994, pelanek_bayesian_2017} & Qual é o estado latente de domínio e o desempenho futuro esperado? & Modelos probabilísticos e sequenciais de \textit{Knowledge Tracing} (como o BKT) sobre trajetórias cronológicas $X = \{(q_1, a_1), \ldots, (q_t, a_t)\}$. & Estimação contínua da probabilidade de domínio $P(L_t)$ por estudante e conceito, atualizada recursivamente a cada nova interação no Cognita. \\
\hline
\textbf{Prescritivo / Orientador} \cite{clow_learning_2012, liao_naviedu_2026} & Que ação pedagógica ou de estudo priorizar a seguir? & Sistemas de apoio à decisão, priorização de estudo para autorregulação do estudante e sinalização de materiais críticos ao docente. & Lista priorizada de conceitos a revisar no Atlas (visão do estudante), ranking de baixa absorção no Analytics e indicação visual de materiais para reforço no Deck Manager. \\
\hline
\end{tabularx}
\legend{Fonte: Elaborado pelos autores com base em \citeonline{clow_learning_2012}, \citeonline{lang_handbook_2017} e \citeonline{liao_naviedu_2026}.}
\end{table}


\subsection{Interoperabilidade de Eventos e Engenharia de Dados Educacionais}
\label{subsec:interoperabilidade-la}

Diferentemente da primeira geração de ambientes virtuais de aprendizagem, na qual todas as interações ocorriam dentro de um único sistema monolítico, ecossistemas educacionais contemporâneos distribuem a experiência pedagógica entre múltiplos serviços especializados, como ferramentas de autoria e publicação de aulas, aplicativos de resposta em tempo real, editores de mapas conceituais, ambientes de exploração pelos estudantes e serviços federados de identidade \cite{lang_handbook_2017, richards_fundamentals_2020}. Nesse cenário distribuído, a viabilidade de análises diagnósticas e preditivas depende diretamente da padronização semântica dos eventos e da engenharia de dados que reconcilia essas fontes heterogêneas.

Para enfrentar a fragmentação de registros entre plataformas, a comunidade de tecnologia educacional estabeleceu padrões de interoperabilidade para fluxos de atividades de aprendizagem, com destaque para a especificação \textit{Experience API} (xAPI) \cite{lang_handbook_2017}. Inspirado na estrutura semântica ator, verbo e objeto (\textit{actor, verb, object}), o modelo de eventos estilo xAPI registra cada experiência de aprendizagem como um fato imutável acompanhado de carimbo de tempo (\textit{timestamp}), contexto instrucional (como sessão, turma ou versão do material) e resultado observável. No Ecossistema CEDIS, essa abordagem é adotada por uma camada dedicada à ingestão de eventos de aprendizagem (Seção~\ref{subsec:arquitetura-dados}).

Sob a perspectiva de arquitetura de dados \cite{kleppmann_designing_2017}, a construção de um \textit{pipeline} de \textit{Learning Analytics} apto a alimentar modelos de \textit{Knowledge Tracing} sobre microsserviços independentes exige o tratamento rigoroso de quatro propriedades:
\begin{itemize}
    \item \textbf{Separação semântica entre trilhas de dados:} distinguir formalmente eventos de telemetria de navegação (como abertura de material ou visualização de resumo), fatos avaliativos de resposta (cargas úteis contendo o desempenho do estudante em uma questão) e vínculos de contexto e identidade, respeitando a finalidade e o ciclo de vida de cada trilha;
    \item \textbf{Reconciliação lógica e idempotência:} em arquiteturas nas quais serviços mantêm bancos de dados isolados, a associação entre o evento de interação, a resposta persistida e o conceito avaliado realiza-se por referências lógicas (\textit{cross-service}). O \textit{pipeline} analítico deve garantir deduplicação idempotente diante de reenvios de rede e validar a integridade referencial antes do treinamento do modelo, evitando que respostas órfãs ou atividades sem vínculo conceitual na Matriz Q contaminem a estimação de parâmetros;
    \item \textbf{Tratamento de heterogeneidade estrutural:} sistemas de resposta em sala de aula frequentemente armazenam diferentes formatos de atividade (múltipla escolha, escalas opinativas do tipo Likert e texto aberto) em campos semiestruturados (JSON). A camada de engenharia de atributos de LA deve validar e filtrar estritamente os itens avaliativos dotados de critério objetivo de correção para derivar a variável binária $a_t \in \{0, 1\}$, auditando em \textit{log} os registros descartados por inconsistência;
    \item \textbf{Coerência temporal e proveniência:} como modelos sequenciais dependem da ordem exata das tentativas do estudante, todos os carimbos de tempo provenientes de serviços distribuídos devem ser normalizados em Tempo Universal Coordenado (UTC), e cada versão de modelo treinado deve registrar o período de extração e a versão dos esquemas de domínio consumidos, assegurando reprodutibilidade científica.
\end{itemize}

\subsection{Ciclo Duplo de \textit{Feedback}, Autorregulação e Explicabilidade}
\label{subsec:ciclo-duplo-la}

O valor pedagógico de um sistema de \textit{Learning Analytics} materializa-se na forma como suas inferências são entregues aos usuários finais para fechar o ciclo de intervenção \cite{clow_learning_2012, verbert_learning_2014}. Em alinhamento com as diretrizes contemporâneas da área, o projeto de intervenções analíticas eficazes apoia-se em um \textbf{ciclo duplo de \textit{feedback}} (\textit{dual feedback loop}), que atende simultaneamente às necessidades metacognitivas do estudante (nível micro) e às demandas de gestão pedagógica do docente (nível macro) \cite{gaevic_learning_2015, liao_naviedu_2026}:

\begin{enumerate}
    \item \textbf{Nível micro (\textit{Learning Analytics} voltado ao estudante e autorregulação da aprendizagem):} Fundamentado na teoria da Aprendizagem Autorregulada (\textit{Self-Regulated Learning}, SRL) de \citeonline{zimmerman_becoming_2002}, o sucesso acadêmico em disciplinas complexas ou no estudo autônomo depende da capacidade do aprendiz de monitorar sua própria compreensão, identificar lacunas e planejar estrategicamente a alocação de seu tempo de estudo. Contudo, estudantes frequentemente apresentam limitações metacognitivas ao tentar diagnosticar a causa raiz de seus erros, sem distinguir se uma falha decorre de um conceito recém-exposto ou de um pré-requisito anterior não consolidado, ou hesitam em expor dúvidas publicamente em sala de aula. Interfaces de \textit{Learning Analytics} voltadas ao estudante (\textit{learner-facing LA}) atuam como andaimes metacognitivos (\textit{scaffolds}) ao externalizar a probabilidade de domínio estimada para cada conceito, ordenar uma lista priorizada de tópicos não dominados para revisão imediata e exibir a evolução histórica do progresso individual, fortalecendo a autoeficácia e a persistência no estudo \cite{verbert_learning_2014, cao_multi-granularity_2025};

    \item \textbf{Nível macro (\textit{Learning Analytics} voltado ao docente e apoio à decisão pedagógica):} Em turmas universitárias numerosas, o professor enfrenta uma barreira de escala para diagnosticar tempestivamente quais conceitos não foram absorvidos pela turma antes da aplicação de avaliações somativas formais \cite{verbert_learning_2014}. Painéis analíticos docentes (\textit{teacher-facing Learning Analytics Dashboards}, LADs) superam as limitações das médias brutas de participação ao consolidar as probabilidades individuais inferidas pelo modelo de \textit{Knowledge Tracing} em indicadores coletivos por turma e curso: proporção estimada de estudantes que atingiram o limiar de domínio em cada conceito, identificação e ranqueamento dos conceitos mais críticos (do menor para o maior índice de absorção), visualização da evolução temporal da absorção ao longo das semanas letivas e indicação direta de quais materiais didáticos (apresentações e atividades) associados aos conceitos deficitários necessitam de reforço em sala.
\end{enumerate}

Um requisito transversal a ambos os polos desse ciclo duplo é a \textbf{transparência e explicabilidade algorítmica}, tema central do subcampo de \textit{Explainable Learning Analytics} (XLA) e Inteligência Artificial Explicável na Educação (\textit{Explainable AI in Education}, XAI-ED) \cite{khosravi_explainable_2022, cao_multi-granularity_2025}. Conforme demonstram \citeonline{khosravi_explainable_2022}, tanto docentes, naturalmente céticos perante pontuações opacas desprovidas de comprovação empírica, quanto estudantes tendem a rejeitar ou interpretar incorretamente diagnósticos preditivos quando estes são apresentados apenas como números isolados de ``caixa-preta''. Para construir confiança e subsidiar decisões pedagógicas responsáveis, a camada de apresentação do \textit{Learning Analytics} deve acompanhar cada estimativa de domínio de três elementos explicativos: (i) o \textbf{grau de confiança} da estimativa (refletindo a quantidade e a consistência das interações observadas); (ii) o \textbf{tamanho da amostra} ($n$ de estudantes considerados) nas visões agregadas de turma; e (iii) os \textbf{motivos pedagógicos da estimativa}, traduzindo os parâmetros interpretáveis do modelo (como o histórico recente de tentativas e acertos, a ocorrência de erros consecutivos ou a influência dos parâmetros de chute e descuido do conceito) em justificativas compreensíveis em linguagem clara.

\subsection{Ética, Privacidade e Governança de Dados Educacionais}
\label{subsec:etica-privacidade-la}

A coleta e a modelagem de dados sobre o comportamento e o estado cognitivo de estudantes suscitam questões éticas e regulatórias profundas, amplamente debatidas na literatura de \textit{Learning Analytics} por autores como \citeonline{slade_learning_2013} e \citeonline{pardo_ethical_2014}. Esses trabalhos estabelecem que sistemas analíticos educacionais não devem operar como mecanismos de vigilância punitiva, rotulação determinística ou exposição pública de vulnerabilidades dos estudantes, mas sim como instrumentos formativos regidos pela transparência, pela autonomia do estudante e pela proteção estrita da privacidade desde a concepção (\textit{Privacy by Design}) \cite{cavoukian_privacy_2009}.

Do ponto de vista técnico e arquitetural, a implementação de \textit{Learning Analytics} com rastreamento de conhecimento sob a vigência da Lei Geral de Proteção de Dados (LGPD) impõe a resolução de uma tensão estrutural entre \textbf{minimização de telemetria} e \textbf{continuidade longitudinal}. Para métricas descritivas de uso geral, a adoção de identificadores efêmeros rotativos (como códigos HMAC recalculados diariamente a partir de atributos truncados de rede e classe de dispositivo), aliada ao respeito ao sinalizador \textit{Do Not Track} (DNT) e à filtragem de tráfego automatizado, é suficiente e desejável, pois impede o rastreamento persistente não autorizado. Todavia, a inferência do estado de conhecimento individual via \textit{Knowledge Tracing} exige, por definição matemática, encadear a sequência histórica de respostas $X = \{(q_1, a_1), \ldots, (q_t, a_t)\}$ de um mesmo aprendiz ao longo de múltiplos dias e sessões de estudo.

Em consonância com os princípios éticos formulados por \citeonline{slade_learning_2013} e \citeonline{pardo_ethical_2014}, o equacionamento dessa tensão requer quatro salvaguardas arquiteturais que orientam diretamente o projeto do módulo Cognita no Ecossistema CEDIS:
\begin{enumerate}
    \item \textbf{Segregação por consentimento explícito (\textit{opt-in}) e pseudonimização:} O processamento longitudinal individual deve ser restrito exclusivamente aos titulares que manifestarem consentimento ativo (\textit{opt-in}) por meio de um serviço de identidade controlado pelo próprio estudante. Mesmo sob consentimento, o motor preditivo deve operar sobre identificadores pseudonimizados, mantendo os modelos analíticos desacoplados de dados civis diretos como nome ou endereço de e-mail;
    \item \textbf{Direito à autodeterminação informativa e expurgo em cascata:} O titular dos dados deve deter controle contínuo sobre seu histórico preditivo, incluindo o direito de revogar o consentimento e solicitar a exclusão definitiva de todas as interações reconciliadas, estimativas de domínio e séries históricas associadas à sua identidade, cabendo à modelagem de dados garantir o expurgo físico atômico dentro de prazos determinísticos;
    \item \textbf{Prevenção de reidentificação indireta via \textit{k}-anonimato:} Em painéis docentes e relatórios agregados por curso, turma ou filtro específico, grupos reduzidos podem revelar indiretamente o desempenho de um estudante individual. A aplicação formal do modelo de \textit{k}-anonimato \cite{sweeney_k-anonymity_2002}, segundo o qual qualquer célula ou recorte analítico só é exibido se contiver registros de ao menos $k$ indivíduos distintos (por exemplo, $k = 5$), previne a exposição individual por cruzamento de filtros;
    \item \textbf{Caráter consultivo e não estigmatizante (\textit{soft guidance}):} As estimativas geradas pelo sistema devem restringir-se à orientação formativa privada ao próprio estudante e à agregação anônima para o docente, vedando-se qualquer ranqueamento nominal entre pares, bloqueio rígido da navegação do estudante (\textit{hard gating}) ou utilização automatizada do índice de domínio como nota somativa oficial.
\end{enumerate}

Dessa forma, o arcabouço conceitual de \textit{Learning Analytics} estabelece a ponte entre a engenharia de software de sistemas distribuídos (Seção~\ref{sec:engenharia-software}) e os algoritmos de aprendizagem de máquina para rastreamento de conhecimento (Seção~\ref{sec:aprendizagem-maquina}), assegurando que a transição da telemetria descritiva para o diagnóstico preditivo ocorra com rigor semântico, utilidade pedagógica comprovada em ciclo duplo e conformidade ética e legal.


\section{Aprendizagem de Máquina}
\label{sec:aprendizagem-maquina}

A Aprendizagem de Máquina (\textit{Machine Learning}, ML) é a subárea da Inteligência Artificial dedicada ao desenvolvimento de programas capazes de melhorar automaticamente seu desempenho a partir da experiência. Na definição clássica de \citeonline{mitchell_machine_1997}, diz-se que um programa aprende com a experiência $E$, em relação a uma classe de tarefas $T$ e a uma medida de desempenho $P$, se o seu desempenho nas tarefas de $T$, medido por $P$, melhora com a experiência $E$. Em vez de codificar explicitamente as regras de decisão, o desenvolvedor especifica um modelo parametrizado e um critério de otimização, deixando que os parâmetros sejam ajustados a partir de dados observados.

Os métodos de aprendizagem de máquina costumam ser organizados em três paradigmas: aprendizagem supervisionada, não supervisionada e por reforço \cite{mitchell_machine_1997}. Interessa a este trabalho a aprendizagem supervisionada, na qual o modelo aprende, a partir de exemplos rotulados, uma função que mapeia entradas a saídas, aplicada a dados sequenciais. A previsão do desempenho de um estudante em uma questão futura constitui um problema de classificação binária (a resposta será correta ou incorreta), cujos rótulos são obtidos diretamente do histórico de respostas registrado nas plataformas educacionais \cite{cui_rebalancing_2025}. O objetivo do treinamento não é memorizar os exemplos observados, mas generalizar para interações ainda não vistas. Por essa razão, os dados são particionados em conjuntos de treino, validação e teste, frequentemente com validação cruzada em $k$ partições (\textit{k-fold cross-validation}), de modo a detectar o sobreajuste (\textit{overfitting}), situação em que o modelo se ajusta excessivamente aos dados de treino e perde capacidade de generalização \cite{shen_assessing_2022, nagatani_augmenting_2019}.

Nas últimas décadas, a aprendizagem profunda (\textit{deep learning}) consolidou-se como a vertente de maior expressividade da área. Fundamentada em redes neurais artificiais com múltiplas camadas, ela aprende representações hierárquicas diretamente dos dados brutos, dispensando boa parte da engenharia manual de atributos \cite{lecun_deep_2015}. Arquiteturas como as Redes Neurais Recorrentes (\textit{Recurrent Neural Networks}, RNN), em especial a variante \textit{Long Short-Term Memory} (LSTM) \cite{hochreiter_long_1997}, e o \textit{Transformer}, baseado em mecanismos de atenção \cite{vaswani_attention_2017}, tornaram-se predominantes na modelagem de sequências. Esse ganho de expressividade, contudo, tem como contrapartida a redução da interpretabilidade: o estado interno de uma rede profunda é um vetor denso cujos componentes, isoladamente, não possuem significado direto para o usuário final \cite{cao_multi-granularity_2025}. Em aplicações educacionais, nas quais docentes e estudantes precisam compreender por que um diagnóstico foi emitido, esse compromisso entre precisão e interpretabilidade é decisivo para a escolha do modelo.

\subsection{Rastreamento de Conhecimento (\textit{Knowledge Tracing})}
\label{subsec:knowledge-tracing}

O Rastreamento de Conhecimento (\textit{Knowledge Tracing}, KT) é a tarefa de modelar a evolução do estado de conhecimento latente de um estudante a partir de suas interações com um sistema de aprendizagem e de prever seu desempenho em interações futuras \cite{corbett_knowledge_1994, shen_assessing_2022}. Trata-se de um componente central de sistemas tutores inteligentes e de plataformas adaptativas, pois a estimativa do domínio de cada conceito fundamenta decisões como a recomendação de exercícios, a priorização de revisões e a identificação precoce de estudantes em risco \cite{liu_enhancing_2026, gao_transformer-based_2025}.

Formalmente, considera-se um conjunto de estudantes, um conjunto de questões e um conjunto de conceitos (\textit{knowledge concepts}, KC), no qual cada questão está associada a um ou mais conceitos por meio de uma matriz de associação questão-conceito, denominada Matriz Q \cite{liu_enhancing_2026}. O histórico de um estudante é representado por uma sequência temporal de interações $X = \{(q_1, a_1), (q_2, a_2), \ldots, (q_t, a_t)\}$, em que $q_i$ é a questão respondida no instante $i$ e $a_i \in \{0, 1\}$ indica se a resposta foi incorreta ou correta. Dado esse histórico, o modelo de KT estima a probabilidade de acerto da próxima questão, $P(a_{t+1} = 1 \mid X, q_{t+1})$, e, de forma subjacente, a probabilidade de domínio de cada conceito \cite{shen_assessing_2022, cui_rebalancing_2025}.

Dois aspectos distinguem o KT de uma classificação convencional. Primeiro, o estado de conhecimento não é diretamente observável: as respostas constituem apenas evidências ruidosas, sujeitas a acertos por chute e a erros por descuido. Segundo, esse estado é dinâmico, pois evolui com a prática e decai com o tempo em razão do esquecimento \cite{nagatani_augmenting_2019}. A literatura propôs diversas famílias de modelos para lidar com esses aspectos, sintetizadas na Tabela~\ref{tab:familias-kt} e detalhadas nas subseções seguintes.

\begin{table}[H]
\caption{Síntese das principais famílias de modelos de \textit{Knowledge Tracing}}
\label{tab:familias-kt}
\centering
\small
\begin{tabularx}{\textwidth}{>{\raggedright\arraybackslash}p{2.6cm}>{\raggedright\arraybackslash}X>{\raggedright\arraybackslash}X>{\raggedright\arraybackslash}X}
\hline
\textbf{Família} & \textbf{Representação do estado de conhecimento} & \textbf{Vantagens} & \textbf{Limitações} \\
\hline
Probabilística (BKT e extensões) \cite{corbett_knowledge_1994, pardos_kt-idem_2011} & Variável latente binária por conceito (domina / não domina), atualizada por um Modelo Oculto de Markov & Parâmetros com significado pedagógico explícito; baixo custo computacional; funciona com poucos dados por conceito & Conceitos tratados de forma independente; ausência de esquecimento no modelo original; menor poder preditivo \\
\hline
Fatorial e psicométrica (IRT, PFA) \cite{embretson_item_2013, pavlik_performance_2009} & Proficiência contínua estimada por função logística a partir de acertos, erros e dificuldade dos itens & Interpretação psicométrica consolidada; incorporação simples de atributos & Desconsidera a ordem completa da sequência de interações; supõe proficiência estável na IRT clássica \\
\hline
Profunda recorrente e de memória (DKT, DKVMN) \cite{piech_deep_2015, zhang_dynamic_2017} & Vetor oculto denso (RNN/LSTM) ou matriz de memória chave-valor & Captura dinâmicas não lineares e dependências temporais longas & Estado latente pouco interpretável; demanda grandes volumes de dados \\
\hline
Baseada em atenção (SAKT, AKT) \cite{pandey_self-attentive_2019, ghosh_context-aware_2020} & Representação contextual ponderada pelas interações históricas mais relevantes & Elevado desempenho preditivo; pesos de atenção oferecem explicação parcial & Custo computacional maior; sensível a dados esparsos \\
\hline
Estrutural e híbrida (grafos, transferência, caminho duplo) \cite{liu_enhancing_2026, liao_naviedu_2026} & Estados por conceito propagados sobre grafos ou hipergrafos de pré-requisitos & Modela relações entre conceitos e generaliza para conceitos pouco observados & Maior complexidade de engenharia e de implantação \\
\hline
\end{tabularx}
\legend{Fonte: Elaborado pelos autores.}
\end{table}


\subsubsection{\textit{Bayesian Knowledge Tracing}}
\label{subsubsec:bkt}

O \textit{Bayesian Knowledge Tracing} (BKT), proposto por \citeonline{corbett_knowledge_1994}, é o modelo fundador da área e permanece como referência pela clareza de sua formulação. O BKT modela o domínio de cada conceito como um Modelo Oculto de Markov (\textit{Hidden Markov Model}, HMM) com dois estados latentes (o estudante domina ou não domina o conceito) e observações binárias correspondentes às respostas corretas ou incorretas \cite{pelanek_bayesian_2017, nagatani_augmenting_2019}. Cada conceito é descrito por quatro parâmetros probabilísticos, todos com interpretação pedagógica direta:

\begin{itemize}
    \item \textbf{Conhecimento inicial ($P(L_0)$):} probabilidade de o estudante já dominar o conceito antes da primeira interação;
    \item \textbf{Aprendizado ($P(T)$):} probabilidade de transição do estado de não domínio para o de domínio após uma oportunidade de prática;
    \item \textbf{Chute ($P(G)$):} probabilidade de o estudante acertar a questão sem dominar o conceito (\textit{guess});
    \item \textbf{Descuido ($P(S)$):} probabilidade de o estudante errar a questão mesmo dominando o conceito (\textit{slip}).
\end{itemize}

A cada resposta observada, o BKT atualiza a crença sobre o domínio do conceito aplicando o teorema de Bayes. Sendo $P(L_t)$ a probabilidade de domínio antes da interação $t$, a probabilidade a posteriori, condicionada a uma resposta correta ou incorreta, é dada respectivamente pelas Equações~\ref{eq:bkt-correta} e~\ref{eq:bkt-incorreta} \cite{corbett_knowledge_1994, pelanek_bayesian_2017}:

\begin{equation}
\label{eq:bkt-correta}
P(L_t \mid a_t = 1) = \frac{P(L_t)\,\bigl(1 - P(S)\bigr)}{P(L_t)\,\bigl(1 - P(S)\bigr) + \bigl(1 - P(L_t)\bigr)\,P(G)}
\end{equation}

\begin{equation}
\label{eq:bkt-incorreta}
P(L_t \mid a_t = 0) = \frac{P(L_t)\,P(S)}{P(L_t)\,P(S) + \bigl(1 - P(L_t)\bigr)\,\bigl(1 - P(G)\bigr)}
\end{equation}

Em seguida, incorpora-se a possibilidade de aprendizado decorrente da prática, obtendo-se a estimativa de domínio para a próxima interação (Equação~\ref{eq:bkt-transicao}) e, a partir dela, a probabilidade de acerto da próxima questão associada ao conceito (Equação~\ref{eq:bkt-predicao}):

\begin{equation}
\label{eq:bkt-transicao}
P(L_{t+1}) = P(L_t \mid a_t) + \bigl(1 - P(L_t \mid a_t)\bigr)\,P(T)
\end{equation}

\begin{equation}
\label{eq:bkt-predicao}
P(a_{t+1} = 1) = P(L_{t+1})\,\bigl(1 - P(S)\bigr) + \bigl(1 - P(L_{t+1})\bigr)\,P(G)
\end{equation}

Os quatro parâmetros de cada conceito são estimados a partir dos dados históricos, tipicamente pelo algoritmo de Maximização da Expectativa (\textit{Expectation-Maximization}, EM) ou por busca em grade, sendo comum a imposição de restrições como $P(G) < 0{,}5$ e $P(S) < 0{,}5$ para evitar soluções degeneradas sem sentido pedagógico \cite{pelanek_bayesian_2017}. O resultado é, para cada estudante e conceito, uma probabilidade de domínio $P(L_t)$ que evolui a cada interação e pode ser apresentada diretamente ao usuário.

Essa formulação confere ao BKT propriedades particularmente relevantes para sistemas educacionais em produção. Os parâmetros são interpretáveis por docentes (por exemplo, um valor elevado de $P(G)$ sinaliza questões com alternativas facilmente adivinháveis), o custo computacional de treinamento e de inferência é baixo e a atualização a cada nova resposta é imediata, o que viabiliza o acompanhamento contínuo do estudante \cite{corbett_knowledge_1994, pelanek_bayesian_2017}. Em contrapartida, o modelo original apresenta limitações conhecidas: representa o conhecimento de forma binária, trata cada conceito de maneira independente, ignorando relações de pré-requisito, desconsidera a dificuldade específica de cada questão e não modela o esquecimento, já que a probabilidade de domínio não decai com o passar do tempo \cite{nagatani_augmenting_2019, cui_rebalancing_2025}.

Diversas extensões foram propostas para mitigar essas limitações sem abandonar a estrutura probabilística. O KT-IDEM (\textit{Item Difficulty Effect Model}) incorpora parâmetros de chute e descuido específicos por questão, introduzindo a dificuldade dos itens no modelo \cite{pardos_kt-idem_2011}. Outras variações acrescentam o efeito do intervalo de tempo entre interações ou do número de tentativas para representar o esquecimento \cite{nagatani_augmenting_2019}. Essas extensões preservam a principal virtude do BKT, que é fornecer uma estimativa de domínio por conceito explicável e de baixo custo, o que o mantém como base adequada para cenários em que a transparência do diagnóstico é tão importante quanto a precisão.

\subsubsection{Modelos Psicométricos e Fatoriais}
\label{subsubsec:modelos-psicometricos}

Em paralelo ao BKT, desenvolveu-se uma família de modelos baseados em funções logísticas, oriunda da psicometria. A Teoria de Resposta ao Item (\textit{Item Response Theory}, IRT) estima a probabilidade de acerto de uma questão $i$ em função da proficiência $\theta$ do estudante e de parâmetros do item, como a discriminação $a_i$ e a dificuldade $b_i$, conforme o modelo logístico de dois parâmetros \cite{embretson_item_2013, cui_rebalancing_2025}:

\begin{equation}
\label{eq:irt-2pl}
P_i(\theta) = \frac{1}{1 + e^{-a_i(\theta - b_i)}}
\end{equation}

A IRT clássica, contudo, supõe que a proficiência do estudante permanece estável durante a avaliação, o que a torna pouco adequada para acompanhar a aprendizagem ao longo da prática diária \cite{shen_assessing_2022}. A \textit{Performance Factors Analysis} (PFA) adapta essa abordagem ao contexto de aprendizagem ao estimar a probabilidade de acerto a partir da quantidade de acertos e erros anteriores do estudante em cada conceito \cite{pavlik_performance_2009}. Embora permita incorporar atributos adicionais com facilidade, a PFA considera apenas interações relativas ao mesmo conceito e não explora a ordem completa da sequência de respostas \cite{nagatani_augmenting_2019}.

\subsubsection{\textit{Deep Knowledge Tracing} e Modelos Profundos}
\label{subsubsec:dkt}

O \textit{Deep Knowledge Tracing} (DKT), proposto por \citeonline{piech_deep_2015}, introduziu a aprendizagem profunda no rastreamento de conhecimento ao modelar a sequência de interações com uma rede LSTM. Diferentemente do BKT, que mantém uma variável binária por conceito, o DKT representa o estado de conhecimento do estudante como o vetor oculto da rede recorrente, denso e de baixa dimensionalidade, a partir do qual se estima a probabilidade de acerto para todos os conceitos simultaneamente \cite{nagatani_augmenting_2019}. Essa representação permite capturar dinâmicas não lineares e relações implícitas entre conceitos, o que resultou em ganhos expressivos de desempenho preditivo em relação aos modelos clássicos \cite{gao_transformer-based_2025}.

A partir do DKT, a literatura evoluiu em diversas direções. O DKVMN (\textit{Dynamic Key-Value Memory Network}) emprega uma matriz estática de chaves para representar os conceitos e uma matriz dinâmica de valores para registrar o domínio do estudante, recuperando parte da interpretabilidade por conceito \cite{zhang_dynamic_2017}. Modelos baseados em atenção, como o SAKT (\textit{Self-Attentive Knowledge Tracing}) \cite{pandey_self-attentive_2019} e o AKT (\textit{Context-Aware Attentive Knowledge Tracing}) \cite{ghosh_context-aware_2020}, inspirados no \textit{Transformer}, ponderam a relevância de cada interação histórica para a previsão atual, sendo que o AKT incorpora ainda parâmetros de dificuldade inspirados no modelo de Rasch.

Trabalhos recentes têm buscado incorporar explicitamente fatores pedagógicos à modelagem. O DIMKT incorpora a dificuldade das questões e dos conceitos e modela o efeito dessa dificuldade antes, durante e após a prática \cite{shen_assessing_2022}. \citeonline{nagatani_augmenting_2019} estendem o DKT com atributos de esquecimento, como o intervalo desde a última interação com o mesmo conceito e o número de tentativas anteriores. \citeonline{cui_rebalancing_2025} demonstram que boa parte dos ganhos de acurácia dos modelos recentes provém de respostas pouco discriminativas, como acertos em questões fáceis, e propõem reponderar o treinamento para privilegiar as respostas que efetivamente diferenciam os níveis de domínio dos estudantes. Na linha estrutural, o MHAKT modela a transferência de conhecimento entre conceitos e áreas por meio de hipergrafos \cite{liu_enhancing_2026}, enquanto o NaviEdu combina um caminho neural e um caminho baseado em grafo de conceitos para detectar mudanças abruptas de desempenho e convertê-las em intervenções direcionadas \cite{liao_naviedu_2026}. Há ainda propostas que acoplam o KT diretamente à recomendação de exercícios \cite{gao_transformer-based_2025} ou que empregam técnicas de explicabilidade \textit{post hoc}, como valores SHAP, para justificar as previsões a docentes \cite{cao_multi-granularity_2025}.

Em comum, esses trabalhos evidenciam que o aumento do desempenho preditivo dos modelos profundos veio acompanhado de maior complexidade arquitetural, maior demanda por dados e menor transparência do estado latente \cite{liu_enhancing_2026, cao_multi-granularity_2025}. Essa tensão é central para a escolha do modelo em um módulo como o Cognita, que precisa entregar diagnósticos compreensíveis a estudantes e docentes a partir de volumes de dados ainda modestos.

\subsection{Avaliação de Modelos de Rastreamento de Conhecimento}
\label{subsec:avaliacao-kt}

Por se tratar de uma tarefa de classificação binária probabilística, os modelos de KT são avaliados com métricas que comparam as probabilidades previstas às respostas efetivamente observadas no conjunto de teste. As mais recorrentes na literatura são \cite{shen_assessing_2022, cui_rebalancing_2025, liao_naviedu_2026}:

\begin{itemize}
    \item \textbf{Área sob a curva ROC (\textit{Area Under the Curve}, AUC):} mede a capacidade do modelo de ordenar corretamente respostas corretas acima das incorretas, independentemente de um limiar de decisão. Varia de 0 a 1, sendo 0,5 o desempenho de um classificador aleatório;
    \item \textbf{Acurácia (ACC):} proporção de respostas cuja correção foi prevista acertadamente, considerando em geral o limiar de 0,5 sobre a probabilidade prevista;
    \item \textbf{Raiz do erro quadrático médio (\textit{Root Mean Squared Error}, RMSE):} mede a distância média entre a probabilidade prevista e o rótulo observado, refletindo a calibração das probabilidades.
\end{itemize}

A interpretação dessas métricas exige cautela. Como a proporção de respostas corretas costuma ser elevada nas bases educacionais, um modelo pode obter acurácia alta apenas acertando respostas previsíveis, sem de fato distinguir estudantes com diferentes níveis de domínio \cite{cui_rebalancing_2025}. Além disso, a divisão dos dados deve ser feita por estudante, e não por interação isolada, para evitar que informações de um mesmo estudante apareçam simultaneamente no treino e no teste \cite{nagatani_augmenting_2019, cao_multi-granularity_2025}. Esses cuidados orientam a definição das métricas de qualidade do Cognita e o monitoramento de cada novo treinamento do modelo.

Em síntese, o rastreamento de conhecimento fornece o arcabouço teórico para transformar o histórico de respostas registrado no Ecossistema CEDIS em estimativas de domínio por conceito. Entre as abordagens apresentadas, o BKT destaca-se pela correspondência direta entre seus parâmetros e o fenômeno pedagógico modelado, pela estimativa explícita da probabilidade de domínio de cada conceito e pelo baixo custo de operação, características alinhadas aos objetivos de transparência e de acompanhamento contínuo do módulo proposto neste trabalho. Por esse motivo, o BKT é adotado como o modelo de rastreamento de conhecimento do Cognita, enquanto os modelos baseados em aprendizagem profunda permanecem na literatura como referência comparativa do estado da arte e possibilidade para extensões futuras do ecossistema.
